\documentclass[11pt]{article}
\usepackage[T1]{fontenc}
\usepackage{lmodern,amsmath,amssymb,amsthm,mathtools,microtype}
\usepackage[margin=1in]{geometry}
\usepackage[colorlinks=true,linkcolor=blue,citecolor=blue,urlcolor=blue]{hyperref}
\numberwithin{equation}{section}
\newtheorem{theorem}{Theorem}[section]
\newtheorem{proposition}[theorem]{Proposition}
\newtheorem{lemma}[theorem]{Lemma}
\theoremstyle{definition}
\newtheorem{definition}[theorem]{Definition}
\theoremstyle{remark}
\newtheorem{remark}[theorem]{Remark}
\newcommand{\R}{\mathbb R}
\newcommand{\Ptwo}{\mathcal P_2(\R^d)}
\newcommand{\dd}{\,\mathrm d}
\newcommand{\Div}{\operatorname{div}}
\newcommand{\esssup}{\operatorname*{ess\,sup}}
\newcommand{\norm}[1]{\lVert #1\rVert}
\newcommand{\Cpsi}{\mathcal C_\psi}
\newcommand{\Dpsi}{\mathcal D_\psi}
\title{Finite propagation in fourth-order aggregation--diffusion}
\author{Kaiwen Jin\thanks{\href{mailto:yui.ui.siki@gmail.com}{\texttt{yui.ui.siki@gmail.com}}}}
\date{}
\begin{document}
\maketitle
\begin{abstract}
We prove finite propagation for a fourth-order aggregation--diffusion equation
with linear mobility and quadratic attraction on $\R^d$.
For every $d\ge2$, every positive attraction coefficient and every compactly
supported nonnegative $H^1$ probability density, there is one global weak
solution whose support remains in a bounded ball on each finite time interval.
The solution has the same $L^\infty H^1\cap L^2H^2$ regularity and integrated
weak formulation as the variational existence theory.
A localized power-entropy estimate is proved on the same minimizing
movements by a bounded-derivative auxiliary flow. The resulting superlinear
inequality for the exterior spacetime tail forces that tail to vanish at a
finite radius, in every finite dimension.
\end{abstract}

\section{Introduction and main theorem}\label{sec:introduction}
We study the degenerate fourth-order equation
\begin{equation}\label{eq:PDE}
 \partial_t\rho=-\Div(\rho\nabla\Delta\rho)-\chi\Delta(\rho^2)
       \quad\hbox{on }(0,\infty)\times\R^d,\qquad \chi>0.
\end{equation}
It describes local aggregation competing with fourth-order diffusion.
The factor $\rho$ in the highest-order flux makes the equation degenerate
at vacuum, while the quadratic term is attractive second-order diffusion.
A basic qualitative question is whether a compact aggregate retains a
bounded edge during its evolution. This does not follow from regularity or
moment bounds alone, and a fourth-order equation has no comparison principle
of the usual second-order kind.

Carrillo, Esposito, Falc\'o and Fern\'andez-Jim\'enez~\cite{CEFF} construct
global probability-valued weak solutions for a family of such equations.
The quadratic attraction in \eqref{eq:PDE} lies in their subcritical range
in every finite dimension. Their construction gives finite second moments
and $H^1,H^2$ control on finite time intervals. In the Introduction of that
work, preservation of compact support is posed as a further question,
supported there by numerical observations. We prove this property for
\eqref{eq:PDE}, retaining their weak formulation and data class.

Finite propagation for thin-film equations is closely connected with
localized entropy estimates. For example, Novick-Cohen and
Shishkov~\cite{NCS} obtain such results for one-dimensional equations with
backwards second-order diffusion. Their interval problem and entropy-solution
construction differ from the whole-space variational weak solution considered
here. The issue is to obtain a sufficiently strong local estimate for that
same weak solution, uniformly in the approximations and without an
upper restriction on the spatial dimension.

Write $\Ptwo$ for probability measures with finite second moment.
Narrow continuity means continuity against every bounded continuous
spatial test. Our main result is the following.

\begin{theorem}\label{thm:main}
Let $d\ge2$ be an integer, $\chi>0$, and let
$\rho_0\in H^1(\R^d)$ be nonnegative, compactly supported, and have
$\int_{\R^d}\rho_0\dd x=1$.
There exists one global narrowly continuous curve
$\rho:[0,\infty)\to\Ptwo$, with $\rho(0)=\rho_0$, represented by
nonnegative probability densities, such that for every finite $T>0$,
\[
 \rho\in L^\infty(0,T;H^1(\R^d))\cap L^2(0,T;H^2(\R^d)).
\]
For every $\phi\in C_c^2(\R^d)$ and every $0\le t_1<t_2<\infty$,
\begin{align}\label{eq:weak}
 \int\phi[\rho(t_2)-\rho(t_1)]\dd x
  =-\int_{t_1}^{t_2}\!\int
   [\rho\Delta\rho\Delta\phi
       +\Delta\rho\nabla\rho\cdot\nabla\phi
                      +\chi\rho^2\Delta\phi]\dd x\dd t.
\end{align}
For each such $T$ there is $R_T<\infty$ for which
\[
 \rho(t,x)=0\quad\hbox{for almost every }
 (t,x)\in(0,T)\times(\R^d\setminus B_{R_T}).
\]
The same curve has these properties for all finite horizons.
\end{theorem}

Here $B_R$ is the Euclidean ball centered at zero. The radius may depend
on the datum, $d,\chi$ and $T$. The new conclusion concerns the support
of a selected global solution; it gives neither uniqueness of the
fourth-order weak equation nor a quantitative growth law for $R_T$.
The theorem retains the integrated weak identity and the
$L^\infty H^1\cap L^2H^2$ bounds of~\cite[Definition~2.1]{CEFF}
at $m=2$.\footnote{Definition~2.1(i) of the pinned preprint also writes
$p=2^*=\infty$ in dimension two. We do not assert that endpoint.
The construction estimates used here are the displayed $H^1/H^2$ bounds
and their consequences for finite spatial exponents.}
The support conclusion is established with precisely the solution
requirements stated above.

The proof uses a power entropy with exponent slightly above one. Its
localized dissipation controls a weighted Hessian and a weighted quartic
gradient. Interpolation turns these into a superlinear estimate for the
exterior spacetime tail. Because that tail tends to zero far away, an
elementary iteration makes it identically zero beyond a finite radius.

At the available weak regularity, direct nonlinear testing does not justify
the entropy calculation. Instead we use the flow-interchange principle
of Matthes, McCann and Savar\'e~\cite{MMS} on the original variational
approximations. A convex density entropy with bounded derivative admits
a spatially weighted metric flow in the sense of Ambrosio, Gigli and
Savar\'e~\cite{AGS}. We establish its energy estimate by renormalization,
positive implicit approximation and comparison in a common energy class.
The resulting local inequality therefore holds for the same original
solution, rather than for a separately chosen entropy solution.

Section~\ref{sec:entropies} introduces these auxiliary entropies and their
spatial estimates. Section~\ref{sec:dissipation} proves the rough-data
auxiliary dissipation. Section~\ref{sec:propagation} obtains the local
inequality and proves the theorem. The two interpolation estimates are
included in Appendix~\ref{app:interpolation}.
Throughout, $H^j=W^{j,2}(\R^d)$ has its inhomogeneous norm;
$\nabla^2u$ is the Hessian with Frobenius norm. Integrals without a domain
are over $\R^d$, or over the explicitly specified periodic cube in the
auxiliary construction. Time suprema are essential suprema. Products
with a negative density power are assigned zero at vacuum; the limiting
arguments below justify each use.

\section{Weighted auxiliary entropies}\label{sec:entropies}

\subsection{The original weak solution and its approximations}
Let
\begin{equation}\label{eq:F}
 F(u)=\frac12\int_{\R^d}|\nabla u|^2\dd x
             -\chi\int_{\R^d}u^2\dd x
\end{equation}
for an $H^1$ probability density, and set $F=+\infty$ otherwise.
The minimizing movement construction of Carrillo, Esposito, Falc\'o and
Fern\'andez-Jim\'enez~\cite[Definition~2.1, Theorem~2.2 and Section~4]{CEFF}
applies to this energy. In their notation $m=2<2+2/d$, for every finite $d$.
Their weak identity has exactly the coefficient $\chi$ in
\eqref{eq:weak}. We recall the estimates of the construction because our
localization will be applied to the same minimizers.

Write $W_2$ for the quadratic Wasserstein distance on $\Ptwo$,
\[
 W_2^2(\mu,\nu)=\inf_{\pi\in\Gamma(\mu,\nu)}
       \int_{\R^d\times\R^d}|x-y|^2\dd\pi(x,y),
\]
where $\Gamma(\mu,\nu)$ is the set of probability measures with marginals
$\mu,\nu$. For $\tau>0$ choose
\begin{equation}\label{eq:JKO}
 u_k\in\operatorname*{argmin}_{v\in\Ptwo}
 \left\{F(v)+\frac{W_2^2(v,u_{k-1})}{2\tau}\right\},
 \qquad u_0=\rho_0,
\end{equation}
and write $u_\tau(t)=u_k$ on $((k-1)\tau,k\tau]$.
For every finite $T$, the source estimates give
\begin{align}
 \sup_{\tau,t\le T}\bigl(\norm{u_\tau(t)}_{H^1}^2
           +\textstyle\int |x|^2u_\tau(t)\dd x\bigr)&\le C_T,
                    \label{eq:JKO-H1}\\
 \int_0^T\norm{u_\tau(t)}_{H^2}^2\dd t&\le C_T,
 &\sum_{k\tau\le T}W_2^2(u_k,u_{k-1})&\le C_T\tau.
                    \label{eq:JKO-H2}
\end{align}
Each $u_k$ has finite entropy and belongs to $H^2$; the latter follows from
the heat-flow estimate in~\cite[Lemma~4.2]{CEFF}. A subsequence converges to
the source weak solution, strongly in spacetime $L^2$. Interpolating this
convergence with \eqref{eq:JKO-H2} gives strong convergence in
$L^2(0,T;H^1)$. A diagonal choice over integer $T$ gives one global curve.
We fix that choice throughout.

Choose the small exponent
\begin{equation}\label{eq:exponents}
 a=\frac1{100(d+1)},\qquad p=1+a,\qquad r=2+3a,\qquad r_1=2+4a.
\end{equation}
For $d>2$ we have $r_1<2d/(d-2)$, while for $d=2$ all these exponents
are finite. The usual Sobolev inequality and mass interpolation therefore
give uniform $L^{r_1}$ bounds. Interpolating with strong $L^2$ convergence
gives strong spacetime $L^r$ convergence along the same subsequence.
Only finite-exponent Sobolev embeddings are used when $d=2$.

\subsection{A bounded derivative and an actual metric flow}
The power entropy $u^p$ has the right localization estimates, but its
unbounded derivative is inconvenient in a spatially weighted transport
flow. We first use a convex approximation with bounded derivative.
For $M\ge1$ and $z\ge0$ define
\begin{equation}\label{eq:scalar}
 q_M(z)=\frac{z^a}{1+(z/M)^{2a}},\qquad
 f_M'(z)=\frac{M^a}{a}\arctan((z/M)^a),\qquad f_M(0)=0.
\end{equation}
Thus $f_M''(z)=q_M(z)/z$ for $z>0$. Put
\begin{equation}\label{eq:pressure}
 P_M(z)=zf_M'(z)-f_M(z)=\int_0^z q_M(y)\dd y,
 \qquad L_M=\frac{\pi M^a}{2a}.
\end{equation}
The elementary bounds needed below are
\begin{gather}
 |zq_M'(z)|\le a q_M(z),\qquad
 0\le f_M(z)\le L_M z,\qquad 0\le f_M'(z)\le L_M,
                                                   \label{eq:scalar-bounds}\\
 zq_M(z)\ge(1-a)P_M(z),\qquad
 zq_M(z)-(1-1/d)P_M(z)\ge(1/d-a)P_M(z),              \label{eq:pressure-bound}\\
 \frac{f_M'(z)^2}{q_M(z)}\le C_a z^a,\qquad
 \frac{f_M'(z)^4}{q_M(z)^3}\le C_a(z^a+z^{3a}).       \label{eq:ratios}
\end{gather}
For \eqref{eq:pressure-bound}, integrate
$q_M(y)\le q_M(z)(z/y)^a$ over $0<y<z$.
For \eqref{eq:ratios}, on $z\le M$ use
$f_M'(z)\le z^a/a$ and $q_M(z)\ge z^a/2$.
On $z\ge M$ use $f_M'\le L_M$ and
$q_M(z)\ge M^{2a}z^{-a}/2$; the remaining factor $M^{-2a}$ is at most one.
Also $q_M\uparrow z^a$ and $f_M\uparrow z^p/(ap)$ as $M\to\infty$.

Let $0\le\zeta\le1$ be smooth and constant outside a compact set, and set
$\psi=\zeta^4$. For $0<\epsilon\le1$ consider
\begin{equation}\label{eq:G}
 G_{M,\epsilon}(u)=\int\psi f_M(u)\dd x
                       +\epsilon\int u\log u\dd x.
\end{equation}
Here $0\log0=0$ and the value is $+\infty$ on singular measures.
The weight will eventually vanish on an inner ball and equal one far away.
The fourth power makes all derivative quotients used below bounded.
For a finite-energy density $u$, the descending metric slope is
\[
 |\partial G_{M,\epsilon}|(u)
 =\limsup_{\substack{v\to u\\v\ne u}}
    \frac{[G_{M,\epsilon}(u)-G_{M,\epsilon}(v)]_+}{W_2(u,v)}.
\]
Here convergence is in $W_2$, and $[\cdot]_+$ denotes the positive part.

\begin{proposition}\label{prop:metric-flow}
For fixed $M,\epsilon,\psi$, the functional \eqref{eq:G} has an
$\lambda_M$-gradient-flow semigroup $S_s$ on its finite-energy domain in
$\Ptwo$, where
\begin{equation}\label{eq:lambda}
 \lambda_M=-L_M\left(\norm{\nabla^2\psi}_\infty
   +(1/d-a)^{-1}\norm{\frac{|\nabla\psi|^2}{\psi}}_\infty\right).
\end{equation}
The quotient at $\psi=0$ is assigned zero. The curve $U_s=S_su$ is
$W_2$-continuous at zero and satisfies the evolution variational inequality
\begin{equation}\label{eq:EVI}
 \frac12\frac{\mathrm d}{\mathrm ds}W_2^2(U_s,v)
      +\frac{\lambda_M}{2}W_2^2(U_s,v)
       \le G_{M,\epsilon}(v)-G_{M,\epsilon}(U_s)
\end{equation}
for almost every $s>0$ and every finite-energy $v$.
For finite initial energy its tangent velocity $v_s$ satisfies
$\partial_sU+\Div(Uv)=0$ and the energy/action identity
\[
 G_{M,\epsilon}(U_s)+\int_0^s\!\int|v_t|^2U_t\dd x\dd t
       =G_{M,\epsilon}(u),\qquad
 \norm{v_s}_{L^2(U_s)}=|\partial G_{M,\epsilon}|(U_s)
\]
with the norm equality holding almost everywhere.
\end{proposition}

\begin{proof}
Abbreviate $G=G_{M,\epsilon}$. We check the hypotheses
of~\cite[Theorem~11.2.1]{AGS}.
For the normalized Gaussian $\gamma$, nonnegativity of relative entropy gives
\begin{equation}\label{eq:entropy-lower}
 \int u\log u\dd x\ge-\frac12\int|x|^2u\dd x-\frac d2\log(2\pi).
\end{equation}
Thus $G$ is bounded below on each second-moment ball and a Gaussian is a
finite-energy test. These are properness and the literal coercivity
condition (11.2.1b) of that theorem.

For completeness the negative entropy tail is uniform on such a ball.
Choose $0<\alpha<2$ and split $\{0<u\le1,\ |x|>R\}$ at
$u=\exp(-|x|^\alpha)$. For large $R$,
\begin{equation}\label{eq:entropy-tail}
 \int_{|x|>R}(u\log u)^-\dd x
 \le R^{\alpha-2}\int|x|^2u\dd x
       +\int_{|x|>R}|x|^\alpha e^{-|x|^\alpha}\dd x.
\end{equation}
Local convex-integral lower semicontinuity consequently extends to the
whole entropy. The nonnegative weighted convex integral in \eqref{eq:G}
is lower semicontinuous as a relaxed integral; the entropy makes its value
infinite on singular measures. One may also use local weak $L^1$
compactness from the entropy bound and weak lower semicontinuity of the
Lipschitz convex functional $u\mapsto\int\psi f_M(u)$;
see~\cite[Corollary~3.9 and Theorem~4.30]{Brezis}.

It remains to verify the precise transport convexity required by the source.
Let $\mu_1=m(x)\dd x$ be a finite-energy base and let $T_2,T_3$ be its
optimal maps to finite-energy endpoints $\mu_2,\mu_3$. The generalized
geodesic under consideration is
\[
 T_t=(1-t)T_2+tT_3,\qquad \mu_t=(T_t)_\#\mu_1,\qquad0\le t\le1,
\]
where the push-forward notation means $\mu_t(E)=\mu_1(T_t^{-1}(E))$.
By~\cite[Lemma~5.5.3 and Proposition~6.2.12]{AGS}, on a common full
base-measure set the maps have positive definite almost-everywhere
Jacobians $A_i=DT_i(x)$, are essentially strictly monotone, and obey the
density-change formula. Put
\[
 A_t=(1-t)A_2+tA_3,\quad D=A_3-A_2=D^{\mathsf T},\quad
 J_t=\det A_t>0.
\]
The interpolated map remains essentially injective. At $y=T_t(x)$ its
density, transported velocity and deformation matrix are respectively
\[
 z_t(y)=m(x)/J_t,\qquad v_t(y)=T_3(x)-T_2(x),\qquad B_t(y)=DA_t^{-1}.
\]
Pulling the entropy back to the base gives the scalar integrand
\[
 K_x(t)=J_t\psi(T_t(x))f_M(m(x)/J_t)
                       +\epsilon m(x)\log(m(x)/J_t).
\]
For regular transport data, the second variation of the weighted energy is
\begin{align}\label{eq:second-variation}
 \int\bigl[\psi\{(zq_M-P_M)(\operatorname{tr}B)^2
                +P_M\operatorname{tr}(B^2)\}
  -2P_M\operatorname{tr}B\,\nabla\psi\cdot v
             +f_M\nabla^2\psi[v,v]\bigr]\dd y,
\end{align}
where $z=z_t(y)$, $v=v_t(y)$, $B=B_t(y)$, and the scalar functions are
evaluated at $z$. The matrix $B$ is similar to the symmetric matrix
$A_t^{-1/2}DA_t^{-1/2}$, so
$\operatorname{tr}(B^2)\ge(\operatorname{tr}B)^2/d$.
By \eqref{eq:pressure-bound} and completion of a square,
\eqref{eq:second-variation} is at least $\lambda_M\int z|v|^2\dd y$.
The entropy has nonnegative second variation.

The same algebra applies pointwise to $K_x$ on the common full set,
before any differentiation is passed through a spatial integral.
It yields $K_x''(t)\ge\lambda_Mm(x)|T_3(x)-T_2(x)|^2$ also for the
nonsmooth transport maps. Integrating first in $t$ and then in $x$ yields
\begin{equation}\label{eq:generalized-convexity}
 G(\mu_t)\le(1-t)G(\mu_2)+tG(\mu_3)
 -\frac{\lambda_M}{2}t(1-t)
         \int|T_3-T_2|^2\dd\mu_1.
\end{equation}
The endpoint energies and induced cost are finite; the negative entropy
part is integrable by \eqref{eq:entropy-tail}. This justifies the integration
also before positive intermediate entropy integrability is known.

The cost in \eqref{eq:generalized-convexity} is the induced three-plan cost.
For negative $\lambda_M$ it is retained in that form.
This is exactly~\cite[Definition~9.2.4 and (11.2.1a)]{AGS}; the distinction
from endpoint cost is explained in its Remark~9.2.6. The generation theorem
itself yields the endpoint-distance inequality \eqref{eq:EVI} and the
energy/action and slope identities (11.2.4), (11.2.9).
\end{proof}

\subsection{A weighted spatial coercivity estimate}
We next compute the decay of $F$ for the auxiliary equation. Its rigorous
application to the rough curve $S_su$ occupies Section~\ref{sec:dissipation}.
For a positive smooth periodic density write $q=q_M(u)$, $f'=f_M'(u)$ and
\begin{equation}\label{eq:CD}
 \Cpsi(u)=\int\psi q|\nabla^2u|^2,\qquad
 \Dpsi(u)=\int\psi\frac q{u^2}|\nabla u|^4.
\end{equation}
The formulas below also define these nonnegative quantities on rough
whole-space densities, with vacuum values zero. Set
\begin{align}
 E_M(u)&=C_d\int u^2\left[
  \left(q+\frac{f'^4}{q^3}\right)\frac{|\nabla\psi|^4}{\psi^3}
       +\frac{f'^2}{q}\frac{|\Delta\psi|^2}{\psi}\right],
                                                    \label{eq:EM}\\
 B_M(u)&=E_M(u)+C_d\chi^2\int\psi q_M(u)u^2
                   +C_d\chi\int u^{2+a}|\Delta\psi|.  \label{eq:BM}
\end{align}
Here and below a sufficiently large $C_d$ can be fixed once; it is
independent of $M,\epsilon$ and the periodic box. These errors have growth
through $u^{2+a}+u^{2+3a}$, by \eqref{eq:ratios}.

\begin{lemma}\label{lem:spatial}
Let $\psi$ be constant near the boundary of a periodic cube and let
$u>0$ be smooth and periodic. For
\begin{equation}\label{eq:aux-flux}
 J(u)=(\epsilon+\psi q_M(u))\nabla u
                          +u f_M'(u)\nabla\psi,
\end{equation}
there is $c_d>0$, independent of $M,\epsilon$ and the cube, such that
\begin{align}\label{eq:spatial-dissipation}
 -\int(-\Delta u-2\chi u)\Div J(u)
 &\ge c_d\bigl(\Cpsi(u)+\Dpsi(u)\bigr)
                +\frac\epsilon2\norm{\Delta u}_2^2\notag\\
 &\hspace{12mm}-B_M(u)-2\chi\epsilon\norm{\nabla u}_2^2.
\end{align}
\end{lemma}

\begin{proof}
Put $g=\nabla u$ and $H=\nabla^2u$. Integrating
$\Div[\psi(q/u)|g|^2g]$ and using $|uq'|\le aq$ gives
\[
 (1-a)\Dpsi\le(2+\sqrt d)\sqrt{\Cpsi\Dpsi}
       +\Dpsi^{3/4}\left(\int qu^2
                        \frac{|\nabla\psi|^4}{\psi^3}\right)^{1/4}.
\]
Consequently
\begin{equation}\label{eq:D-by-C}
 \Dpsi\le\frac{4(2+\sqrt d)^2}{(1-a)^2}\Cpsi
              +C_d\int qu^2\frac{|\nabla\psi|^4}{\psi^3}.
\end{equation}
A second integration by parts gives
\begin{align}\label{eq:Hessian-identity}
 \int\psi\{q(\Delta u)^2+q'|g|^2\Delta u\}
 =\Cpsi+\int\psi q'H[g,g]
       +\int q\nabla\psi\cdot Hg-\int q\Delta u\nabla\psi\cdot g.
\end{align}
Expanding $\int\Delta u\Div J(u)$ and using
\eqref{eq:Hessian-identity}, its non-heat terms are
\begin{equation}\label{eq:fourth-spatial}
 \Cpsi+\int\psi q'H[g,g]+
 \int q\nabla\psi\cdot Hg
 +\int(f'+q)\Delta u\nabla\psi\cdot g
 +\int uf'\Delta\psi\Delta u.
\end{equation}
The first error is at least $-a\sqrt{\Cpsi\Dpsi}$.
Our choice of $a$ satisfies
$2a(2+\sqrt d)/(1-a)<1/4$. Combining \eqref{eq:D-by-C}
with Young's inequality retains a fixed positive multiple of
$\Cpsi+\Dpsi$.
The next two terms in \eqref{eq:fourth-spatial} are first absorbed into
$\Cpsi$ and then into $\Dpsi$. Their remainders are bounded by
$C_du^2(q+f'^4/q^3)|\nabla\psi|^4/\psi^3$ after integration.
The final term leaves $C_du^2 f'^2|\Delta\psi|^2/(q\psi)$.
These are exactly the terms in \eqref{eq:EM}.

The attraction contributes
\[
 -2\chi\int\psi q|g|^2-2\chi\int uf'\nabla\psi\cdot g.
\]
The first term is absorbed into a fraction of $\Dpsi$ with remainder
$C_d\chi^2\int\psi qu^2$. For the second use
$H_M(z)=\int_0^z y f_M'(y)\dd y\le z^{2+a}/[a(2+a)]$.
It equals $2\chi\int H_M(u)\Delta\psi$ and is bounded below by the
last error in \eqref{eq:BM}. The heat contribution is
$\epsilon\norm{\Delta u}_2^2-2\chi\epsilon\norm{\nabla u}_2^2$.
Decreasing $c_d$ once and retaining half the positive heat term proves
\eqref{eq:spatial-dissipation}.
\end{proof}

\section{Dissipation along the auxiliary flow}\label{sec:dissipation}

Fix $M,\epsilon,\psi$ for this section and abbreviate $G=G_{M,\epsilon}$.
The preceding spatial calculation must be established along the actual
flow generated by Proposition~\ref{prop:metric-flow}, starting from a rough
density. We prove the following statement.

\begin{proposition}\label{prop:auxiliary-energy}
Let $\bar u\in H^1(\R^d)$ be a nonnegative probability density in $\Ptwo$,
and let $U_s=S_s\bar u$. For every finite $S$, this curve belongs to
$L^\infty(0,S;H^1)\cap L^2(0,S;H^2)$ and satisfies, for every $s\ge0$,
\begin{align}\label{eq:aux-energy}
 F(U_s)+\int_0^s\left[c_d(\Cpsi(U_t)+\Dpsi(U_t))
                     +\frac\epsilon2\norm{\Delta U_t}_2^2\right]\dd t
 \le F(\bar u)+\int_0^s
       [B_M(U_t)+2\chi\epsilon\norm{\nabla U_t}_2^2]\dd t.
\end{align}
The constants and errors are those of Lemma~\ref{lem:spatial}.
Moreover $U_s\to\bar u$ strongly in $H^1$ and $L^r$ as $s\downarrow0$, and
\begin{align}\label{eq:initial-dissipation}
 \liminf_{s\downarrow0}\frac{F(\bar u)-F(S_s\bar u)}s
 &\ge c_d(\Cpsi(\bar u)+\Dpsi(\bar u))
       +\frac\epsilon2\norm{\Delta\bar u}_2^2\notag\\
 &\hspace{8mm}-B_M(\bar u)-2\chi\epsilon\norm{\nabla\bar u}_2^2.
\end{align}
The positive terms on the right may initially have value $+\infty$.
\end{proposition}

We first obtain a common energy class for the actual metric flow. We then
construct positive periodic approximations for which the spatial estimate
is legitimate. A direct comparison in $H^{-1}$ identifies their limit with
the metric flow and permits the energy inequality to pass to it.

\subsection{The flux and a renormalized identity}
Set
\begin{equation}\label{eq:A}
 \mathcal A(x,z)=\epsilon z+\psi(x)P_M(z),\qquad
 \kappa(x,z)=\epsilon+\psi(x)q_M(z).
\end{equation}
Thus $\epsilon\le\partial_z\mathcal A=\kappa\le\Lambda$,
where $\Lambda=\epsilon+\norm{q_M}_\infty$.

\begin{lemma}\label{lem:actual-flux}
The curve $U=S_s\bar u$ satisfies
\begin{equation}\label{eq:auxiliary-PDE}
 \partial_s U=\Div J,\qquad
 J=\kappa(x,U)\nabla U+Uf_M'(U)\nabla\psi
   =\nabla\mathcal A(x,U)+f_M(U)\nabla\psi.
\end{equation}
On every finite interval, both $J/\sqrt U$ and $\nabla U/\sqrt U$ belong
to spacetime $L^2$, with value zero at vacuum.
\end{lemma}

\begin{proof}
At almost every positive time the metric-flow velocity $v$ is tangent,
$\eta=-v$ is a subgradient, and
$\norm{\eta}_{L^2(U)}=|\partial G|(U)=:\sigma$.
These are the conclusions of~\cite[Definition~11.1.1 and
Theorem~11.2.1]{AGS}. At an absolutely continuous measure, the reduced
subgradient agrees with the ordinary one by its Remark~10.3.3.
For clarity, at an absolutely continuous finite-slope density the ordinary
subgradient condition is~\cite[Definition~10.1.1, (10.1.3)]{AGS}
\[
 G(\nu)-G(U)\ge\int\eta(x)\cdot(X(x)-x)U(x)\dd x
                          +o(W_2(U,\nu))\quad\text{as }\nu\to U,
\]
where $X$ is the optimal transport map from $U$ to the finite-energy
measure $\nu$, and convergence is in $W_2$. Only optimal perturbations
will be used in this inequality.
The tangent space is the $L^2(U)$ closure of compact smooth gradients;
the cylindrical formulation in~\cite[Definition~8.4.1]{AGS} is equivalent
in $\R^d$ by multiplying a bounded cylindrical function by cutoffs tending
to one. The gradient error tends to zero in $L^2(U)$.

For $\xi\in C_c^\infty(\R^d;\R^d)$, direct change of variables gives
\begin{equation}\label{eq:transport-variation}
 \mathcal L(\xi):=\left.\frac{\mathrm d}{\mathrm dh}
       G((\mathrm{Id}+h\xi)_\#U)\right|_{h=0}
 =\int f_M(U)\nabla\psi\cdot\xi
        -\int[\psi P_M(U)+\epsilon U]\Div\xi.
\end{equation}
Both signs of small $h$ give diffeomorphisms. The derivative is dominated
by an integrable multiple of $U$, using $f_M\le L_MU$,
$P_M\le L_MU$ and the bounded logarithmic Jacobian for the entropy.
The metric descending-slope definition and
$W_2(U,(\mathrm{Id}+h\xi)_\#U)\le|h|\norm{\xi}_{L^2(U)}$ give
\[
 |\mathcal L(\xi)|\le\sigma\norm{\xi}_{L^2(U)}.
\]
Choose the sign decreasing the differentiable energy to obtain this bound.
If the norm is zero, the pushforward is the identity $U$-almost everywhere.
By density and Riesz representation,
$\mathcal L(\xi)=\int z\cdot\xi\,U$ for some
$z\in L^2(U)$ with $\norm z\le\sigma$.

For $\xi=\nabla\phi$, $\phi\in C_c^\infty$, the map
$\mathrm{Id}+h\nabla\phi$ is optimal for both sufficiently small signs:
it is the gradient of the convex function $|x|^2/2+h\phi$.
Fenchel's inequality proves optimality, since equality holds on its graph
and both the potential and its conjugate differ from a quadratic by a
bounded function. Their marginal integrals are finite in $\Ptwo$.
The ordinary subgradient inequality can therefore be applied for those
gradient perturbations and yields
$\mathcal L(\nabla\phi)=\int\eta\cdot\nabla\phi\,U$.
The tangent projection of $z$ is consequently $\eta$.
Since $\norm\eta=\sigma\ge\norm z$, orthogonal norm decomposition forces
$z=\eta$. Thus, for all vector tests, \eqref{eq:transport-variation} says
\[
 U\eta=\nabla[\epsilon U+\psi P_M(U)]+f_M(U)\nabla\psi.
\]
The velocity is $v=-\eta$, so its continuity equation has exactly the sign
in \eqref{eq:auxiliary-PDE}. The energy/action identity gives
$J/\sqrt U\in L^2$. Unit mass then gives $J\in L^1$ in spacetime.

The increasing inverse $\mathcal T(x,V)$ of $\mathcal A(x,z)$ has
$|\mathcal T_V|\le\epsilon^{-1}$ and
$|\mathcal T_x|\le C|V|$.
Since $\mathcal A(x,U)$ and its distributional gradient are $L^1$,
the spatial chain rule, first on truncated $V$ and then by these integrable
bounds, gives $U\in W^{1,1}$ in space and the first flux formula in
\eqref{eq:auxiliary-PDE}. At vacuum both $J$ and $\nabla U$ vanish
almost everywhere. Finally
\begin{equation}\label{eq:Fisher}
 \int_0^S\!\int\frac{|\nabla U|^2}{U}
 \le\frac2{\epsilon^2}
       \left[\int_0^S\!\int\frac{|J|^2}{U}
                    +S L_M^2\norm{\nabla\psi}_\infty^2\right]<\infty.
\end{equation}
\end{proof}

\begin{lemma}\label{lem:renormalization}
If $\ell\ge2$ and $\bar u\in L^\ell$, then for every $0\le s\le t\le S$,
\begin{align}\label{eq:Lp-identity}
 \norm{U_t}_\ell^\ell
 +\ell(\ell-1)\int_s^t\!\int U^{\ell-2}\kappa(x,U)|\nabla U|^2
 &=\norm{U_s}_\ell^\ell+\int_s^t\!\int H_\ell(U)\Delta\psi,
\end{align}
where $H_\ell(z)=\ell(\ell-1)\int_0^z y^{\ell-1}f_M'(y)\dd y$.
In particular
\begin{equation}\label{eq:actual-Lp}
 \norm{U_t}_\ell^\ell\le
 e^{(\ell-1)L_M\norm{\Delta\psi}_\infty t}\norm{\bar u}_\ell^\ell.
\end{equation}
Thus $U\in L^\infty L^2\cap L^2H^1$, its flux is $L^2$, and
$\partial_s U\in L^2H^{-1}$ on finite intervals.
\end{lemma}

\begin{proof}
The time chain rule is proved before removing a density truncation. Define
\[
 \beta_n(z)=\begin{cases}
 z^\ell,&0\le z\le n,\\
 \ell n^{\ell-1}z-(\ell-1)n^\ell,&z>n.
 \end{cases}
\]
Its first derivative and $z\beta_n''(z)$ are bounded for fixed $n$;
the latter has only one jump, at $n$. Put
$H_n(z)=\int_0^z y f_M'(y)\beta_n''(y)\dd y$.
We claim the exact identity
\begin{equation}\label{eq:truncated-chain}
 \int\beta_n(U_t)+\int_s^t\!\int\beta_n''(U)\kappa(x,U)|\nabla U|^2
 =\int\beta_n(U_s)+\int_s^t\!\int H_n(U)\Delta\psi.
\end{equation}

Here are the convolution and trace details. Let $\eta_\delta$ be a
nonnegative spatial mollifier and $U_\delta=\eta_\delta*U$.
For any $V$ vanishing at vacuum with $|V|^2/U\in L^1$, weighted
Cauchy--Schwarz gives
\begin{equation}\label{eq:weighted-convolution}
 \frac{|\eta_\delta*V|^2}{U_\delta}
          \le\eta_\delta*\left(\frac{|V|^2}{U}\right).
\end{equation}
Choose a common subsequence of approximate identities converging almost
everywhere for $U$, $V$ and $|V|^2/U$. On $\{U>0\}$ the quotient vector converges
to $V/\sqrt U$. On vacuum its square is bounded by a quantity tending
to zero in \eqref{eq:weighted-convolution}. Fatou and the integrated upper
bound show convergence of its $L^2$ norm. Weak compactness identifies its
weak limit, hence the convergence is strong. Applying this to $J$ and
$\nabla U$, allowed by Lemma~\ref{lem:actual-flux}, gives
\begin{equation}\label{eq:weighted-strong}
 \frac{J_\delta}{\sqrt{U_\delta}}\longrightarrow\frac J{\sqrt U},
 \qquad
 \frac{\nabla U_\delta}{\sqrt{U_\delta}}
                \longrightarrow\frac{\nabla U}{\sqrt U}
 \quad\hbox{strongly in spacetime }L^2.
\end{equation}

For fixed $\delta$,
$\partial_s U_\delta=(\nabla\eta_\delta)*J\in L^1$ in spacetime.
Thus $U_\delta$ has an absolutely continuous $L^1$ representative.
It agrees at every time with $\eta_\delta*U_s$: both are distributionally
continuous in time, the latter by narrow continuity. Its initial value is
exactly $\eta_\delta*\bar u$.
One may now convolve in time on an interior interval, apply the scalar
chain rule for $\beta_n$, and pass back. The derivative difference is
$L^1$, and the remaining factor converges against that fixed derivative
because $\beta_n'$ is bounded and continuous. The Lipschitz bound of
$\beta_n$ passes endpoints. Absolute continuity extends the identity to
$s=0$ and $t=S$.

Multiply the smoothed equation by
$\chi_R\beta_n'(U_\delta)$, where $\chi_R$ is a compact cutoff with
$|\nabla\chi_R|\le C/R$, and integrate in space and time. The boundary
term is at most
$C R^{-1}\norm{\beta_n'}_\infty\int_0^S\int|J|$ and vanishes.
The other integrand is
\[
 [U_\delta\beta_n''(U_\delta)]
       \frac{\nabla U_\delta}{\sqrt{U_\delta}}
                  \cdot\frac{J_\delta}{\sqrt{U_\delta}}.
\]
The product of the weighted vectors converges strongly in $L^1$ by
\eqref{eq:weighted-strong}. Away from $U=n$ its bounded multiplier
converges pointwise. On $U=n$, $\nabla U=0$ almost everywhere by Sobolev
locality, so the limiting product is zero and the jump leaves no error.
This proves strong $L^1$ convergence of the integrand.
At each endpoint, spatial convolution converges in $L^1$ for its actual
density. Consequently
\[
 \int\beta_n(U_t)-\int\beta_n(U_s)
             =-\int_s^t\!\int\beta_n''(U)\nabla U\cdot J
\]
for every endpoint pair, on the same convolution subsequence. Substitution
of \eqref{eq:auxiliary-PDE} and
$\nabla H_n(U)=Uf_M'(U)\beta_n''(U)\nabla U$ gives
\eqref{eq:truncated-chain} by integration against compactly supported
$\nabla\psi$.

The estimate uniform in $n$ is
\[
 0\le H_n(z)\le L_M[z\beta_n'(z)-\beta_n(z)]
                 \le L_M(\ell-1)\beta_n(z).
\]
Below $n$ the pressure is $(\ell-1)z^\ell$; above $n$ it is
$(\ell-1)n^\ell$, while $\beta_n(z)\ge n^\ell$.
Gronwall first gives
$\int\beta_n(U_t)\le e^{(\ell-1)L_M\norm{\Delta\psi}_\infty t}
\int\bar u^\ell$.
The functions $\beta_n$ increase to $z^\ell$, so monotone convergence
proves \eqref{eq:actual-Lp} and spacetime $L^\ell$ integrability.
Now $H_n(U)\le L_M(\ell-1)U^\ell$ is an integrable majorant for the
drift. The positive diffusion integrand increases at positive densities.
Dominated and monotone convergence in \eqref{eq:truncated-chain} prove
\eqref{eq:Lp-identity}. Taking $\ell=2$ proves the stated energy class;
bounded coefficients in the flux then give its $L^2$ bound and the
$H^{-1}$ time derivative. Taking $\ell=r_1$ is justified by
\eqref{eq:exponents} and the initial $H^1$ norm.
\end{proof}

\subsection{Comparison in the energy class}
We will construct a second solution of \eqref{eq:auxiliary-PDE} with an
energy inequality. The following elementary comparison identifies it.

\begin{lemma}\label{lem:comparison}
Two nonnegative narrowly continuous probability-density solutions of
\[
 \partial_s U=\Delta\mathcal A(x,U)+\Div(f_M(U)\nabla\psi)
\]
with the same initial measure and in $L^\infty L^2\cap L^2H^1$ on finite
intervals coincide.
\end{lemma}

\begin{proof}
For two solutions $U,W$ put
$w=U-W$, $C=\mathcal A(x,U)-\mathcal A(x,W)$, and
$D=(f_M(U)-f_M(W))\nabla\psi$.
The coefficient bounds give $C\in L^2H^1$, $D\in L^2$, and
$w_s=\Delta C+\Div D\in L^2H^{-1}$.
Thus $w$ is an $H^1$ time curve in the Hilbert space $H^{-1}$.
Its continuous representative has zero initial trace: narrow continuity
identifies that trace in distributions.

Define the inhomogeneous resolvent $R=(I-\Delta)^{-1}$ variationally by
\[
 Z=Rw\in H^1,\qquad
 \int Z\phi+\nabla Z\cdot\nabla\phi=\langle w,\phi\rangle
 \quad(\phi\in H^1).
\]
This is Riesz representation for the $H^1$ inner product.
For $E=\norm w_{H^{-1}}^2=\norm Z_2^2+\norm{\nabla Z}_2^2$,
the Hilbert time chain rule and the resolvent equation tested with $C$ give
\[
 \frac{E'}2=-\int\nabla C\cdot\nabla Z-\int D\cdot\nabla Z
             =\int CZ-\int Cw-\int D\cdot\nabla Z.
\]
Monotonicity and Lipschitzness yield
$\int Cw\ge\epsilon\norm w_2^2$,
$\norm C_2\le\Lambda\norm w_2$ and
$\norm D_2\le L_M\norm{\nabla\psi}_\infty\norm w_2$.
Therefore
\[
 \frac{E'}2\le-\frac\epsilon2\norm w_2^2
       +\frac{(\Lambda+L_M\norm{\nabla\psi}_\infty)^2}{2\epsilon}E.
\]
Since $E(0)=0$, Gronwall gives $w=0$. Narrow continuity identifies the
representatives at every time.
\end{proof}

\subsection{Positive implicit approximations}
Let $Q_L=(-L,L)^d$ be periodic, with $L\ge1$ large enough that $\psi$ is
constant near its boundary. The derivatives of its periodic restriction
have the original bounds. Extend $P_M,q_M,f_M$ by zero for negative
arguments, so $\mathcal A(x,z)=\epsilon z$ there. The increasing inverse
$\mathcal T$ has slopes between $\Lambda^{-1}$ and $\epsilon^{-1}$.
Also $B(z):=zf_M'(z)$ is globally Lipschitz with constant at most
$L_M+\norm{q_M}_\infty$.

Given a nonnegative periodic previous density $w\in L^2(Q_L)$, an implicit step is
\begin{equation}\label{eq:implicit-step}
 u-h\Delta\mathcal A(x,u)-h\Div(f_M(u)\nabla\psi)=w.
\end{equation}
Here $h$ is an auxiliary construction step. It is separate from the
original minimizing movement step $\tau$ in \eqref{eq:JKO}.

\begin{lemma}\label{lem:implicit-existence}
Let $w\in L^2(Q_L)$ be nonnegative and periodic.
For all sufficiently small $h>0$, depending only on the fixed
$M,\epsilon,\psi$, equation \eqref{eq:implicit-step} has a unique
periodic $H^1$ solution. It is nonnegative and preserves mass.
If $0<k_{\rm old}\le w\le K_{\rm old}<\infty$, its solution satisfies
\begin{equation}\label{eq:barriers}
 \frac{k_{\rm old}}{1+hC_0}\le u\le
       \frac{K_{\rm old}}{1-hC_0},\qquad
 C_0=L_M\norm{\Delta\psi}_\infty,
\end{equation}
and belongs to $H^2(Q_L)$. The restriction on $h$ is independent of
$L,k_{\rm old},K_{\rm old}$.
\end{lemma}

\begin{proof}
Use the variable $V=\mathcal A(x,u)$ and the Hilbert norm
$\norm V_h^2=\norm V_2^2+h\norm{\nabla V}_2^2$ on periodic $H^1$.
The transformed operator is
\[
 \langle K_h(V),\phi\rangle
   =\int\mathcal T(x,V)\phi+h\nabla V\cdot\nabla\phi
                     +h f_M(\mathcal T(x,V))\nabla\psi\cdot\nabla\phi.
\]
For $D=V-W$ its monotonicity estimate is
\begin{align*}
 \langle K_h(V)-K_h(W),D\rangle
 &\ge\left[\frac1\Lambda
       -\frac{hL_M^2\norm{\nabla\psi}_\infty^2}{2\epsilon^2}\right]
                   \norm D_2^2+\frac h2\norm{\nabla D}_2^2\\
 &\ge c\norm D_h^2,\qquad
 c=\min\{1/(2\Lambda),1/2\}>0
\end{align*}
for small $h$. Its dual Lipschitz constant in this norm is at most
$K=\epsilon^{-1}+1+\sqrt h L_M\norm{\nabla\psi}_\infty/\epsilon$.
Let $R_h=I-h\Delta$ be the Riesz map for this Hilbert norm.
The map
$V\mapsto V-\delta R_h^{-1}(K_h(V)-w)$ is a contraction for
$0<\delta<2c/K^2$, as follows by expanding its squared norm.
The contraction theorem gives the solution of $K_h(V)=w$.
The inverse chain gives $u\in H^1$, because $\mathcal T_V$ is bounded and
$|\mathcal T_x(x,V)|\le C|V|$.

Testing with $\min(V,0)$ eliminates the drift on the negative set,
where $\mathcal T=V/\epsilon$. The resulting nonnegative sum is at most
zero, so $V,u\ge0$. Testing with one gives mass conservation.
The step can also be written as
\[
 u-h\Div[\kappa(x,u)\nabla u+B(u)\nabla\psi]=w.
\]
The upper constant in \eqref{eq:barriers} is a supersolution and the lower
constant is a subsolution. Subtract the supersolution and test with
$z=(u-K)_+$. With $K_B=L_M+\norm{q_M}_\infty$,
\[
 \norm z_2^2+h\epsilon\norm{\nabla z}_2^2
       \le hK_B\norm{\nabla\psi}_\infty\norm z_2\norm{\nabla z}_2.
\]
Small $h$ and Young's inequality force $z=0$. The lower barrier follows
by the negative-part test. Thus on $jh\le S$ the successive densities
stay between $k_{\rm initial}e^{-C_0S}$ and
$K_{\rm initial}e^{2C_0S}$, provided $hC_0\le1/2$.

The transformed equation gives
\[
 -\Delta V=(w-u)/h+\Div(f_M(u)\nabla\psi)\in L^2.
\]
Periodic Fourier expansion gives $V\in H^2$, and
$0\le V\le\Lambda u$ makes it bounded. For any bounded periodic $H^2$
function $V$,
\begin{equation}\label{eq:bounded-H2}
 \norm{\nabla V}_4^2
             \le(2+\sqrt d)\norm V_\infty\norm{\nabla^2V}_2.
\end{equation}
For smooth $V$ this follows by integrating
$\Div(V|\nabla V|^2\nabla V)$ and applying Cauchy--Schwarz.
Periodic convolution, preserving the $L^\infty$ bound, proves the general
case. On the positive bounded range of $u$ the inverse $\mathcal T$ has
bounded derivatives through order two. Its second derivative is now
legitimate since $\nabla^2V$ and $|\nabla V|^2$ are $L^2$.
Hence $u\in H^2$, and \eqref{eq:bounded-H2} also gives $\nabla u\in L^4$.
Positive periodic convolution then extends every identity in
Lemma~\ref{lem:spatial} to these steps. In particular, the gradients converge
in $L^4$: apply \eqref{eq:bounded-H2} to the convolution difference, whose
$L^\infty$ norm is bounded and whose Hessian tends to zero in $L^2$.
\end{proof}

The estimates needed to remove the boxes and positive floors do not use
their maximum and minimum. For each fixed $\ell\ge2$, convexity and
testing the step with $\ell u^{\ell-1}$ give
\begin{align}\label{eq:discrete-Lp}
 \norm u_\ell^\ell-\norm w_\ell^\ell
  +h\ell(\ell-1)\int u^{\ell-2}\kappa(x,u)|\nabla u|^2
 &\le h\int H_\ell(u)\Delta\psi\notag\\
 &\le h(\ell-1)L_M\norm{\Delta\psi}_\infty\norm u_\ell^\ell.
\end{align}
For small $h$, iteration gives the uniform bound
$\sup_{jh\le S}\norm{u_j}_\ell^\ell\le
e^{2(\ell-1)L_M\norm{\Delta\psi}_\infty S}\norm{u_0}_\ell^\ell$.
In particular, $\ell=2$ gives
\begin{equation}\label{eq:discrete-flux-bound}
 \sum_{jh\le S}h\norm{\nabla u_j}_2^2\le C_S,
 \qquad \sum_{jh\le S}h\norm{J(u_j)}_2^2\le C_S.
\end{equation}

For moment tightness use the genuinely periodic test
\[
 Z_L(x)=\sum_{i=1}^d(2L/\pi)^2\sin^2(\pi x_i/(2L)).
\]
On $Q_L$ it satisfies
$(4/\pi^2)|x|^2\le Z_L(x)\le|x|^2$,
$|\Delta Z_L|\le2d$ and $|\nabla Z_L|\le2|x|$.
Equation \eqref{eq:implicit-step}, $\mathcal A(x,u)\le\Lambda u$ and
mass one imply
\begin{equation}\label{eq:periodic-moment}
 \int Z_L(u_j-u_{j-1})
 \le h\left[2d\Lambda+
       2L_M\sup_{\operatorname{supp}\nabla\psi}|x||\nabla\psi(x)|\right].
\end{equation}
The physical second moments of the densities restricted to $Q_L$ are
therefore uniformly bounded. Convexity of the entropy and the test
$1+\log u_j$ give
\begin{equation}\label{eq:discrete-entropy}
 H(u_j)-H(u_{j-1})+h\epsilon\int\frac{|\nabla u_j|^2}{u_j}
 \le h\int f_M(u_j)\Delta\psi\le hL_M\norm{\Delta\psi}_\infty.
\end{equation}
Their zero extensions, considered only as densities and not Sobolev
extensions, obey \eqref{eq:entropy-lower}. Thus \eqref{eq:periodic-moment}
and bounded initial entropy give uniform Fisher information summed in
time, and also
\begin{equation}\label{eq:discrete-action}
 \sum_{jh\le S}h\int\frac{|J(u_j)|^2}{u_j}\le C_S.
\end{equation}

We retain the time-discretization defect in the nonconvex energy $F$.
Let $z=u-w=h\Div J(u)$, which has mean zero. Periodic Fourier duality gives
\[
 \norm z_2^2\le\norm z_{\dot H^{-1}}\norm{\nabla z}_2,
 \qquad \norm z_{\dot H^{-1}}\le h\norm{J(u)}_2.
\]
Here the homogeneous norm uses frequencies $\pi n/L$, with the zero mode
omitted. The exact quadratic identity yields
\begin{align}\label{eq:quadratic-defect}
 F(u)-F(w)
 &=\langle-\Delta u-2\chi u,z\rangle
            -\frac12\norm{\nabla z}_2^2+\chi\norm z_2^2\notag\\
 &\le h\langle-\Delta u-2\chi u,\Div J(u)\rangle
                              +\chi^2h^2\norm{J(u)}_2^2.
\end{align}
The pairing is an $L^2$ pairing: the step equation gives
$\Div J=(u-w)/h\in L^2$ and $u\in H^2$.
Lemma~\ref{lem:spatial} now gives
\begin{align}\label{eq:discrete-F}
 F(u)+h\left[c_d(\Cpsi(u)+\Dpsi(u))
                      +\frac\epsilon2\norm{\Delta u}_2^2\right]
 \le F(w)+h[B_M(u)+2\chi\epsilon\norm{\nabla u}_2^2]
                        +\chi^2h^2\norm{J(u)}_2^2.
\end{align}
The sum of the final terms is $O(h)$ by
\eqref{eq:discrete-flux-bound}. The density errors are bounded using
\eqref{eq:discrete-Lp} with $\ell=r_1$, and the gradient error by
\eqref{eq:discrete-flux-bound}. The periodic mass-one coercivity estimate
\begin{equation}\label{eq:periodic-coercivity}
 F(u)\ge\tfrac14\norm{\nabla u}_2^2-C_{d,\chi},\qquad L\ge1,
\end{equation}
proved in Lemma~\ref{lem:periodic-coercivity}, is uniform in the box.
Summing \eqref{eq:discrete-F} gives uniform $\sup H^1$,
$\sum h\norm{u_j}_{H^2}^2$ and weighted positive dissipation bounds.
All constants in these construction bounds may depend on the fixed
$M,\epsilon,\psi,S$; none depends on the floor or box size.

\subsection{Whole-space convergence and the initial dissipation}
Choose $L_j\to\infty$ and smooth compactly supported probability densities
$\rho_j$, supported well inside $Q_{L_j}$, converging to $\bar u$ in
$H^1,L^{r_1}$ and second moment. Smooth expanding cutoffs, normalization,
and a sufficiently fine centered convolution give these approximations.
Set $\alpha_j=L_j^{-4}$ and
\[
 u_{0,j}=(1-\alpha_j)\rho_j+\frac{\alpha_j}{(2L_j)^d}.
\]
The floors have second moment $O(\alpha_jL_j^2)\to0$, zero gradient and
vanishing $L^2$ norm. Thus $F(u_{0,j})\to F(\bar u)$; initial moments,
$L^{r_1}$ norms and entropy bounds are uniform.
Choose $h_j\to0$ small enough for the preceding restrictions and construct
steps through $S+1$. Write $u_j$ for their piecewise constant density,
$J_j$ for its flux, and $v_j$ for the piecewise affine interpolation.
As fields or measures only, extend them by zero outside the physical cube.
The estimates above are
\begin{gather}
 \sup_{j,t}\left[\norm{u_j(t)}_{H^1(Q_{L_j})}
      +\norm{u_j(t)}_{r_1}+\int|x|^2u_j(t)\dd x\right]\le C,\notag\\
 \int_0^{S+1}\!\int_{Q_{L_j}}
    \left[|\nabla^2u_j|^2+|J_j|^2+\frac{|J_j|^2}{u_j}\right]\dd x\dd t
       \le C.\label{eq:approx-bounds}
\end{gather}
Both interpolations have mass one and the same uniform moment and spatial
norm bounds, while $\partial_t v_j=\Div J_j$ on each step.
For a compact smooth test $\phi$, eventually supported inside the cube,
\[
 \left|\int\phi[v_j(t)-v_j(s)]\right|
 \le\norm{\nabla\phi}_\infty\sqrt{|t-s|}
                 \left(\int_s^t\!\int\frac{|J_j|^2}{u_j}\right)^{1/2}.
\]
Tightness and a diagonal Arzel\`a argument for a countable determining
family of tests give one narrowly continuous probability curve $W_t$ with
initial density $\bar u$ and bounded second moment.

Here the local strong convergence can be seen directly. Fix a smooth
compact cutoff $\eta$ and place $\eta v_j$ on one fixed periodic box.
It is bounded in $L^2H^1$, and
$\partial_t(\eta v_j)=\Div(\eta J_j)-\nabla\eta\cdot J_j$
is bounded in $L^2H^{-1}$. The spacetime $L^2$ norm squared of Fourier
modes above frequency $K$ is at most $CK^{-2}$. Each retained coefficient
is bounded in time $H^1$ and hence has a uniform one-half H\"older modulus.
Finite-dimensional compactness followed by $K\to\infty$ proves local
strong $L^2$ convergence. On each step
$\norm{v_j-u_j}_{H^{-1}_{\rm loc}}\le Ch_j\norm{J_j}_2$.
The Fourier interpolation inequality between $H^{-1}$ and $H^1$ shows
that $u_j$ has the same local strong $L^2$ limit. The local $L^2H^2$
bound then gives strong local gradients. A subsequence converges almost
everywhere in density and gradient. At every fixed time the uniform
$H^1,L^{r_1}$ bounds and the prescribed narrow limit give a density $W_t$
with those same bounds.

To pass the energy errors, local convergence is insufficient. Put
$\vartheta=(r_1-2)/(r_1-1)>0$. Mass interpolation and the moment bound give
\begin{equation}\label{eq:density-tail}
 \int_{|x|>R}u_j^2
 \le\left(\int_{|x|>R}u_j\right)^\vartheta
         \left(\int u_j^{r_1}\right)^{1/(r_1-1)}
 \le CR^{-2\vartheta}.
\end{equation}
The same holds for $W$, so the densities converge strongly in global
spacetime $L^2$ and, by interpolation, in $L^r$.
For gradients choose $\theta_R=0$ on $B_{R/2}$ and $\theta_R=1$ outside
$B_R$, with $|\nabla\theta_R|\le C/R$. When $L_j>2R$ it is constant
near the periodic boundary. Periodic integration by parts gives
\begin{equation}\label{eq:gradient-tail}
 \int\theta_R^2|\nabla u_j|^2
 \le2\norm{\theta_Ru_j}_2\norm{\theta_R\Delta u_j}_2
                +CR^{-2}\int_{|x|>R/2}u_j^2.
\end{equation}
After time integration, \eqref{eq:approx-bounds} and
\eqref{eq:density-tail} make this tail tend to zero uniformly.
The Hessian fields, extended as $L^2$ fields, have a global weak limit
identified with $\nabla^2W$ by interior compact tests. Hence $W\in L^2H^2$
and the analogous gradient-tail bound holds for it. Combining local
gradient convergence with these tails gives global strong gradient
convergence. This reasoning does not differentiate a zero extension.

Bounded continuity of $q_M$ and Lipschitzness of $B(z)=zf_M'(z)$ now give
\[
 J_j\longrightarrow\kappa(x,W)\nabla W+Wf_M'(W)\nabla\psi
       \quad\hbox{strongly in global spacetime }L^2.
\]
Pass the compact-test equation to obtain \eqref{eq:auxiliary-PDE} for $W$,
with its initial measure. Both $W$ and the actual $U$ belong to the class
of Lemma~\ref{lem:comparison}; thus $W=U$ at every time.

Sum \eqref{eq:discrete-F} and pass to this limit. At almost every endpoint
$t$, the densities converge in global $L^2$ and the gradients are weakly
compact, so $F(U_t)\le\liminf F(u_j(t))$. For the positive spacetime terms,
$\sqrt{q_M(u_j)}$ is bounded and converges almost everywhere; multiplying
an $L^2$ test therefore converges strongly, which gives weighted Hessian
lower semicontinuity. Fatou passes the quartic term on $\{U>0\}$, and
its target value at vacuum is zero. The heat term is weakly lower
semicontinuous.

Each density integrand in $B_M$ is continuous at zero, has growth bounded
by $C(z^{2+a}+z^{2+3a})$, and converges in $L^1$ by the global strong
$L^2,L^r$ convergence. On bounded density ranges its derivative is bounded
by $C_Kz$; the remaining tails are uniformly integrable by strong $L^r$.
The positive gradient error on the right converges by the global strong
gradient result \eqref{eq:gradient-tail}. The summed defect is $O(h_j)$.
An endpoint shift by at most $h_j$ tends to zero after the $L^1$ convergence
of the right-hand integrands. These facts prove \eqref{eq:aux-energy}
at almost every endpoint.

It holds at every endpoint. Along any sequence $t_k\to t$, the uniform
$H^1,L^{r_1}$ and moment bounds give global strong $L^2$ compactness;
narrow continuity identifies its limit with $U_t$. Weak gradient
compactness gives $F(U_t)\le\liminf F(U_{t_k})$. Choose $t_k$ from the
full-measure endpoint set and use continuity of the indefinite $L^1$
integrals. This extends \eqref{eq:aux-energy} to all $t$, including zero.

Near zero its right-hand error is $O(t)$ by the uniform $H^1,L^{r_1}$
bounds. Thus $\limsup_{t\downarrow0}F(U_t)\le F(\bar u)$.
The preceding compactness and narrow continuity give strong $L^2$
convergence and the reverse lower semicontinuity, hence $F(U_t)\to F(\bar u)$.
The gradient norms then converge, proving strong $H^1$ at zero.
Interpolation with the uniform $L^{r_1}$ bound gives strong $L^r$.

Finally set
\[
 \mathfrak D(w)=c_d(\Cpsi(w)+\Dpsi(w))
                         +\frac\epsilon2\norm{\Delta w}_2^2,
 \qquad
 \mathfrak E(w)=B_M(w)+2\chi\epsilon\norm{\nabla w}_2^2.
\]
The errors $\mathfrak E(U_t)$ converge to $\mathfrak E(\bar u)$ at zero.
The extended nonnegative functional $\mathfrak D$ is lower semicontinuous
under this strong $H^1$ convergence: a finite liminf controls $H^2$ since
$\epsilon>0$, the weighted Hessian passes weakly, and the quartic term
passes by Fatou and almost-everywhere gradients.
Rearranging \eqref{eq:aux-energy} gives
\[
 \frac{F(\bar u)-F(U_t)}t
       \ge\frac1t\int_0^t[\mathfrak D(U_s)-\mathfrak E(U_s)]\dd s.
\]
Lower semicontinuity and error continuity prove
\eqref{eq:initial-dissipation}, also when $\mathfrak D(\bar u)=+\infty$.
This completes the proof of Proposition~\ref{prop:auxiliary-energy}.
\section{Localized entropy and finite propagation}\label{sec:propagation}

We return to the original minimizing movements \eqref{eq:JKO}. All
auxiliary steps and periodic boxes used in Section~\ref{sec:dissipation}
have already been removed.

\begin{proposition}\label{prop:local-entropy}
The original global weak curve can be selected so that, for each cutoff
$\zeta$ in any prescribed countable family of smooth functions taking
values in $[0,1]$ and constant outside compact sets, and almost every $t>0$,
\begin{align}\label{eq:local-entropy}
 \int\zeta^4u(t)^p
 &+c_d\int_0^t\!\int\zeta^4
       [u^a|\nabla^2u|^2+u^{a-2}|\nabla u|^4]\notag\\
 &\le\int\zeta^4\rho_0^p
 +C_d\int_0^t\!\int u^{p+1}
       [|\nabla\zeta|^4+\zeta^2|\nabla^2\zeta|^2+\chi^2\zeta^4].
\end{align}
The positive constants depend only on $d$ through the fixed exponent
\eqref{eq:exponents}; they are independent of the cutoff radius, width,
time and approximation parameters.
\end{proposition}

\begin{proof}
Use the flow-interchange inequality
of~\cite[Theorem~3.2, (3.21)]{MMS}. For the actual $\lambda_M$-flow of $G$,
with
$D^GF(v)=\limsup_{s\downarrow0}[F(v)-F(S_sv)]/s$, it reads
\[
 G(u_k)+\tau D^GF(u_k)
       +\frac{\lambda_M}{2}W_2^2(u_k,u_{k-1})\le G(u_{k-1}).
\]
Every $u_k$ is an $H^1$ density in $\Ptwo$ with finite entropy, so all
domains match. Proposition~\ref{prop:auxiliary-energy} bounds this
dissipation from below. Its possibly infinite positive terms must then
be finite at a minimizer, since the displayed upper bound is finite.
Summation, dropping the positive $\epsilon H^2$ term, gives
\begin{align}\label{eq:interchange}
 G(u_N)+c_d\tau\sum_{k=1}^N[\Cpsi(u_k)+\Dpsi(u_k)]
 &\le G(\rho_0)
   +\tau\sum_{k=1}^N[B_M(u_k)+2\chi\epsilon\norm{\nabla u_k}_2^2]
       \notag\\
 &\quad+\frac{|\lambda_M|}{2}\sum_{k=1}^NW_2^2(u_k,u_{k-1}).
\end{align}
At fixed $M,\epsilon$ the final term is $O_M(\tau)$ by
\eqref{eq:JKO-H2}.

First let $\tau\to0$ along the original global subsequence. Strong
spacetime $H^1,L^r$ convergence passes the density and gradient errors.
For the positive terms use the bounded continuous weight $q_M$ with weak
Hessians, and almost-everywhere gradient convergence with Fatou for the
quartic term. At almost every endpoint the lower semicontinuity of $G$
and the common moment bound apply. These bounds also give a uniform lower
entropy bound. Letting $\epsilon\to0$ next therefore removes the entropy
and heat error and leaves the weighted $f_M$ inequality for the same
original curve.

Finally $M\to\infty$. The functions $q_M,f_M$ increase to $u^a,u^p/(ap)$.
The error coefficients in \eqref{eq:EM} converge pointwise to constants
times $u^a$ and are dominated using \eqref{eq:ratios} by
$C_d(u^{2+a}+u^{2+3a})$. Monotone convergence on positive terms and
dominated convergence on errors are thus legitimate. Multiplication by
$ap$ gives \eqref{eq:local-entropy}. Indeed
\[
 \frac{|\nabla\psi|^4}{\psi^3}=256|\nabla\zeta|^4,\qquad
 \frac{|\Delta\psi|^2}{\psi}
       \le C_d(|\nabla\zeta|^4+\zeta^2|\nabla^2\zeta|^2),
\]
and Young's inequality bounds $\chi|\Delta\psi|$ by the same cutoff
terms plus $C_d\chi^2\zeta^4$.
Countable cutoff families, integer horizons, integer $M$ and a sequence
$\epsilon\downarrow0$ suffice, so a common refinement preserves one
global original curve and all its source weak-solution properties.
\end{proof}

Fix in Proposition~\ref{prop:local-entropy} the countable family of radial
cutoffs with inner radius $s=m+1-2^{-j}$ and width $\delta=2^{-j-1}$,
for integers $m\ge1$ and $j\ge0$. Each is zero on $B_s$, one outside
$B_{s+\delta}$, and satisfies $|\nabla\zeta|\le C_d\delta^{-1}$ and
$|\nabla^2\zeta|\le C_d\delta^{-2}$. These bounds are uniform for $s\ge1$.
Only these pairs $(s,\delta)$ are needed below.

Fix $T>0$ and choose $R_0\ge1$ containing the essential support of
$\rho_0$. For $s\ge R_0$ let
\begin{equation}\label{eq:tail}
 A(s)=\int_0^T\!\int_{|x|>s}u^{p+1}\dd x\dd t.
\end{equation}
It is finite and tends to zero as $s\to\infty$, by the uniform $H^1$
bound and the strict subcritical exponent. For any of the fixed pairs
$(s,\delta)$ with $s\ge R_0$, use the corresponding cutoff $\zeta$.
The initial term in \eqref{eq:local-entropy} vanishes, giving
\begin{align}\label{eq:tail-energy}
 \esssup_{0<t<T}\int\zeta^4u^p
 +\int_0^T\!\int\zeta^4
       [u^a|\nabla^2u|^2+u^{a-2}|\nabla u|^4]
 \le C_d(\delta^{-4}+\chi^2)A(s).
\end{align}
For the time integral, let admissible endpoints tend to $T$.

\begin{lemma}\label{lem:power-chain}
Put
\begin{equation}\label{eq:power-choices}
 b=\frac{p+1}{2}=1+\frac a2,\qquad
 q=\frac{2p}{p+1},\qquad k=\frac{2(p+1)}p>2,\qquad
 w=\zeta^k u^b.
\end{equation}
Then $w\in L^2(0,T;H^2(\R^d))$ and
\begin{equation}\label{eq:w-bound}
 \esssup_{t<T}\norm{w(t)}_q^q
             +\int_0^T\norm{\nabla^2w(t)}_2^2\dd t
        \le C_d(\delta^{-4}+\chi^2)A(s).
\end{equation}
\end{lemma}

\begin{proof}
The identities $kq=4$, $bq=p$ give $|w|^q=\zeta^4u^p$.
We supply the second-order chain justification because no unweighted
$L^4$ gradient embedding is available in high dimension.
Regularize the power at zero by
$\Phi_\eta(z)=(z+\eta)^b-\eta^b$.
Multiply its derivative by a smooth upper cutoff $\theta(z/K)$, equal to
one below $K$ and zero above $2K$, and integrate from zero, with $K\ge\eta$.
The resulting $\Phi_{\eta,K}$ has bounded first and second derivatives,
vanishes at zero, and satisfies, with constants independent of $\eta,K$,
\begin{gather*}
 |\Phi'|^2\le C_d(z^a+\eta^a),\qquad
 |\Phi''|^2\le C_d z^{a-2}\quad(z>0),\\
 |\Phi(z)|^2\le C_d(z^{2+a}+\eta^a z^2).
\end{gather*}
On $K<z<2K$ the added second derivative has the same bound since $z$
is comparable to $K$.
Writing $\sigma=\zeta^k$, the distributional identity is
\begin{align}\label{eq:power-Hessian}
 \nabla^2[\sigma\Phi(u)]
 &=\sigma\Phi'(u)\nabla^2u
   +\sigma\Phi''(u)\nabla u\otimes\nabla u\notag\\
 &\quad+2\operatorname{sym}[\nabla\sigma\otimes\Phi'(u)\nabla u]
                       +\Phi(u)\nabla^2\sigma.
\end{align}
Initially this is an $L^1_{\rm loc}$ chain/product identity: the truncated
first derivative is bounded in spatial $W^{1,2}_{\rm loc}$, and its
product with $\nabla u\in H^1$ is in $W^{1,1}_{\rm loc}$.
The estimates now show that its right side is $L^2$.

Since $2k\ge4$, the first two squared terms are bounded in spacetime by
\[
 C_d\int_0^T\!\int\zeta^4
       [u^a|\nabla^2u|^2+u^{a-2}|\nabla u|^4]
            +C_d\eta^a\int_0^T\!\int|\nabla^2u|^2.
\]
For the cross term use
\[
 \zeta^{2k-2}|\nabla\zeta|^2u^a|\nabla u|^2
 \le c\zeta^4u^{a-2}|\nabla u|^4
          +C_{d,c}\zeta^{4k-8}|\nabla\zeta|^4u^{a+2}.
\]
Here $4k-8>0$. The $\eta^a$ part is bounded by a vanishing multiple
of $\int_0^T\norm{\nabla u}_2^2$.
Finally
\[
 |\nabla^2\sigma|^2\le C_d[
     \zeta^{2k-4}|\nabla\zeta|^4
                   +\zeta^{2k-2}|\nabla^2\zeta|^2].
\]
This is supported in the transition annulus and bounded by
$C_d\delta^{-4}$. Its contribution is at most
$C_d\delta^{-4}A(s)$ plus a vanishing $\eta^a$ multiple of
$\int_0^T\norm u_2^2$.
First let $K\to\infty$ and then $\eta\downarrow0$.
The powers converge in spacetime $L^2$ and the Hessians are weakly
compact in $L^2$; all $\eta^a$ terms vanish by the original $H^1,H^2$
bounds. Thus the Hessian estimate passes to $w$.
Global $L^2$ and Hessian $L^2$ imply inhomogeneous $H^2$ by derivative
interpolation. Combining with \eqref{eq:tail-energy} proves
\eqref{eq:w-bound}.
\end{proof}

\begin{proof}[Proof of Theorem~\ref{thm:main}]
For the exponent $q$ in \eqref{eq:power-choices},
Lemma~\ref{lem:GN} gives
\[
 \norm w_2\le C_d\norm{\nabla^2w}_2^\theta\norm w_q^{1-\theta},
 \qquad \theta=\frac d{d+4p}.
\]
Integrate its square in time, use H\"older and \eqref{eq:w-bound}, and note
that $w=u^b$ outside $B_{s+\delta}$. We obtain
\begin{equation}\label{eq:tail-recurrence}
 A(s+\delta)\le C_d T^\gamma
      (\delta^{-4}+\chi^2)^\beta A(s)^\beta,
 \qquad
 \gamma=\frac{4p}{d+4p},\qquad
 \beta=1+\frac4{d+4p}>1.
\end{equation}
Indeed $\gamma=1-\theta$ and
$\theta+2(1-\theta)/q=\beta$.

To see exactly how this forces a zero tail, choose an integer
$s_*\ge R_0$ and put
$s_j=s_*+1-2^{-j}$, $\delta_j=2^{-j-1}$, $y_j=A(s_j)$ for $j\ge0$.
Equation \eqref{eq:tail-recurrence} implies
\[
 y_{j+1}\le KB^j y_j^\beta,\qquad
 K=C_dT^\gamma(1+\chi^2)^\beta2^{4\beta},\quad B=2^{4\beta}.
\]
If $y_0\le K^{-1/(\beta-1)}B^{-1/(\beta-1)^2}$, induction gives
\[
 y_j\le K^{-1/(\beta-1)}
          B^{-j/(\beta-1)-1/(\beta-1)^2}\longrightarrow0.
\]
Such a finite integer $s_*$ exists because $A(s)\to0$.
Monotonicity then gives $A(s_*+1)=0$.
Nonnegativity implies that $u=0$ almost everywhere on
$(0,T)\times(\R^d\setminus B_{s_*+1})$.

This uses only the countable family fixed before the tail estimate.
Taking integer time horizons in the preceding argument gives the support
property for the same original global curve on all those horizons;
any finite $T$ is contained in a larger integer horizon. Its narrow continuity, probability and moment properties,
regularity and exact weak identity are the unchanged CEFF conclusions
from Section~\ref{sec:entropies}. This proves every clause of
Theorem~\ref{thm:main}.
\end{proof}
\appendix
\section{Two interpolation estimates}\label{app:interpolation}

\begin{lemma}\label{lem:GN}
If $1\le q<2$ and $w\in H^2(\R^d)\cap L^q(\R^d)$, then
\[
 \norm w_2\le C_{d,q}\norm{\nabla^2w}_2^\theta\norm w_q^{1-\theta},
 \qquad
 \theta=\frac{d(1/q-1/2)}{2+d(1/q-1/2)}.
\]
\end{lemma}
\begin{proof}
Choose an even nonnegative smooth unit kernel $\eta_h$ and put
$s=d(1/q-1/2)$, $L=\norm w_q$, $D=\norm{\nabla^2w}_2$.
Convolution is $L^q$-contractive and
$\norm{\eta_h*w}_\infty\le C h^{-d/q}L$.
Interpolating these two bounds gives
$\norm{\eta_h*w}_2\le C h^{-s}L$.
The first moment of the kernel is zero, so Taylor's integral remainder
and translation continuity in $L^2$ give
$\norm{w-\eta_h*w}_2\le C h^2D$.
This is first proved for smooth $w$ and passed by $H^2$ approximation.
Thus $\norm w_2\le C(h^{-s}L+h^2D)$ for every $h>0$.
If either norm is zero, the conclusion follows directly, letting
$h\to\infty$ when $D=0$. Otherwise choose
$h=(L/D)^{1/(s+2)}$. This proves the formula; for
$q=2p/(p+1)$ it gives $\theta=d/(d+4p)$ in every finite dimension.
\end{proof}

\begin{lemma}\label{lem:periodic-coercivity}
For a nonnegative periodic $H^1(Q_L)$ density of mass one and $L\ge1$,
\[
 \frac12\norm{\nabla u}_2^2-\chi\norm u_2^2
                      \ge\frac14\norm{\nabla u}_2^2-C_{d,\chi}.
\]
\end{lemma}
\begin{proof}
Let $c_n=(2L)^{-d}\int_{Q_L}u(x)e^{-i\pi n\cdot x/L}\dd x$.
Mass one gives $|c_n|\le(2L)^{-d}$.
Splitting Parseval's sum at physical frequency $K>0$, a lattice count
and the gradient norm give
\[
 \norm u_2^2\le C_d(K^d+L^{-d})+K^{-2}\norm{\nabla u}_2^2.
\]
With $E=\norm{\nabla u}_2^2>0$ choose $K=E^{1/(d+2)}$ to obtain
$\norm u_2^2\le C_dE^{d/(d+2)}+C_dL^{-d}$.
The constant-density case $E=0$ has the same bound.
Since $d/(d+2)<1$, scalar Young's inequality proves the assertion
uniformly for $L\ge1$.
\end{proof}

\section*{AI-use statement}
This manuscript was developed with substantial AI assistance in mathematical
exploration, proof drafting, source checking and exposition. Two fresh AI
contexts provided reader and referee-style feedback on a frozen draft,
followed by author revision. These AI checks do not constitute independent
human review or formal verification. Independent mathematical assessment
of the argument is requested.
\begin{thebibliography}{9}
\bibitem{AGS}
L. Ambrosio, N. Gigli and G. Savar\'e,
\emph{Gradient Flows in Metric Spaces and in the Space of Probability Measures},
first edition, Birkh\"auser, 2005.
\href{https://www2.stat.duke.edu/~sayan/ambrosio.pdf}{Institutional copy}.

\bibitem{Brezis}
H. Brezis, \emph{Functional Analysis, Sobolev Spaces and Partial Differential
Equations}, Springer, 2011.
\href{https://doi.org/10.1007/978-0-387-70914-7}{doi:10.1007/978-0-387-70914-7}.

\bibitem{CEFF}
J. A. Carrillo, A. Esposito, C. Falc\'o and A. Fern\'andez-Jim\'enez,
Competing effects in fourth-order aggregation--diffusion equations,
\emph{Proc. Lond. Math. Soc.} \textbf{129} (2024), e12623.
\href{https://doi.org/10.1112/plms.12623}{doi:10.1112/plms.12623}.
Theorem and section numbering used here is from
\href{https://arxiv.org/abs/2307.14706v1}{arXiv:2307.14706v1}.

\bibitem{MMS}
D. Matthes, R. J. McCann and G. Savar\'e,
A family of nonlinear fourth order equations of gradient flow type,
\emph{Comm. Partial Differential Equations} \textbf{34} (2009), 1352--1397.
\href{https://www.math.toronto.edu/mccann/papers/MattesMcCannSavare.pdf}{Author copy}.

\bibitem{NCS}
A. Novick-Cohen and A. Shishkov,
The thin film equation with backwards second order diffusion,
\emph{Interfaces Free Bound.} \textbf{12} (2010), 463--496.
\href{https://doi.org/10.4171/IFB/242}{doi:10.4171/IFB/242}.
\end{thebibliography}
\end{document}
